
\section*{9. Quotient sheaves}
\phantomsection\label{sec:sga3-VIB-9}%
\addcontentsline{toc}{section}{9. Quotient sheaves}

The present section is essentially limited to recalling, in the
particular case of homogeneous spaces of groups, well-known general
facts concerning passage to the quotient by flat equivalence relations
(cf.\ Expos\'e~IV).

\medskip
\phantomsection\label{def:sga3-VIB-9.1}%
\noindent\textbf{Definition 9.1.---}
\emph{Given a monomorphism \(u:G'\to G\) of \(S\)-groups, we
write \(G/G'\) \textup{(respectively, \(G'\backslash G\))} and call it
the \textup{right \textup{(respectively, left)} quotient sheaf of
\(G\) by \(G'\)}, meaning the quotient sheaf, for the
\textup{(fpqc)} topology, of \(G\) by the equivalence relation defined
by the monomorphism}
\phantomsection\label{form:sga3-VIB-9.1-equivalence-relations}%
\[
\begin{aligned}
  G\times_S G'
    &\xrightarrow{\ \delta\circ(\operatorname{id}_G\times u)\ }
      G\times_S G,\\[-1mm]
  &\hspace{7em}\textup{\emph{\textup{(respectively,}}}\\[-1mm]
  G'\times_S G
    &\xrightarrow{\ \gamma\circ(u\times\operatorname{id}_G)\ }
      G\times_S G\textup{\emph{\textup{),}}}
\end{aligned}
\]
\emph{where \(\delta\) \textup{(respectively, \(\gamma\))} denotes
the automorphism of \(G\times_SG\) defined, for
\(g,h\in G(T)\), by}
\[
  (g,h)\longmapsto(g,gh)
  \qquad
  \bigl(\textup{respectively, }(h,g)\longmapsto(hg,g)\bigr).
\]

\medskip
\phantomsection\label{prop:sga3-VIB-9.2}%
\noindent\textbf{Proposition 9.2.---}
\emph{Let \(u:G'\to G\) be a monomorphism of \(S\)-groups.
Suppose that \(G/G'\) is representable by an \(S\)-scheme \(G''\).
Then:}
\begin{enumerate}
\item[\textup{(i)}]\phantomsection\label{concl:sga3-VIB-9.2-i}
  \emph{The canonical morphism \(p:G\to G''\) is a covering for
  the \textup{(fpqc)} topology.}

\item[\textup{(ii)}]\phantomsection\label{concl:sga3-VIB-9.2-ii}
  \emph{If \(\varepsilon''=p\circ\varepsilon\)
  \textup{(this morphism is called the unit section of \(G''\))},
  the following diagrams are cartesian:}
  \begin{center}
    \phantomsection\label{diag:sga3-VIB-9.2-cartesian}%
    \begin{tabular}{c@{\qquad}c}
      \begin{tikzcd}[row sep=large,column sep=large]
        G\times_SG'
          \arrow[r,"{\mu\circ(\operatorname{id}_G\times u)}"]
          \arrow[d,"\operatorname{pr}_1"'] &
        G \arrow[d,"p"] \\
        G \arrow[r,"p"'] &
        G''
      \end{tikzcd},
      &
      \begin{tikzcd}[row sep=large,column sep=large]
        G' \arrow[r,"u"] \arrow[d,"\pi'"'] &
        G \arrow[d,"p"] \\
        S \arrow[r,"\varepsilon''"'] &
        G''
      \end{tikzcd}.
    \end{tabular}
  \end{center}
  \emph{In particular, \(u\) is an immersion.}

\item[\textup{(iii)}]\phantomsection\label{concl:sga3-VIB-9.2-iii}
  \emph{There exists on \(G''\) a unique structure of an
  \(S\)-scheme with left operator group \(G\) such that \(p\) is a
  morphism of \(S\)-schemes with operator group \(G\).}

\item[\textup{(iv)}]\phantomsection\label{concl:sga3-VIB-9.2-iv}
  \emph{If, in addition, \(G'\) is assumed invariant in \(G\),
  there exists on \(G''\) a unique \(S\)-group structure such that
  \(p\) is a morphism of \(S\)-groups.}

\item[\textup{(v)}]\phantomsection\label{concl:sga3-VIB-9.2-v}
  \emph{Let \(S_0\) be an \(S\)-scheme, and put
  \(G_0=G\times_SS_0\) and \(G'_0=G'\times_SS_0\). Then
  \(G_0/G'_0\) is representable by
  \(G''_0=G''\times_SS_0\).}%

\item[\textup{(vi)}]\phantomsection\label{concl:sga3-VIB-9.2-vi}
  \emph{Let \(\mathscr P\) be a property of an \(S\)-morphism.
  Suppose that \(\mathscr P\) is stable under base change. If
  \(p:G\to G''\) has property \(\mathscr P\), then the structural
  morphism \(\pi':G'\to S\) has property \(\mathscr P\) as well.}

\item[\textup{(vii)}]\phantomsection\label{concl:sga3-VIB-9.2-vii}
  \emph{Let \(\mathscr P\) be a property of an \(S\)-morphism.
  Suppose that \(\mathscr P\) is local for the \textup{(fpqc)}
  topology \textup{(cf.\
  \hyperref[par:sga3-VIB-2.0]{2.0} and
  \hyperref[rem:sga3-VIB-2.1.2]{2.1.2})}. Then \(p:G\to G''\)
  has property \(\mathscr P\) if and only if
  \(\pi':G'\to S\) has property \(\mathscr P\).}%
  \begingroup
  \renewcommand{\thefootnote}{\fnsymbol{footnote}}%
  \footnote[1]{\label{note:sga3-VIB-9.2-vii-base}
  Translator's note: the French source prints \(S'\) as the target of
  \(\pi'\) in this clause. No \(S'\) is defined here, and clause
  \textup{\SGARefLink{concl:sga3-VIB-9.2-vi}{(vi)}} identifies \(\pi'\) as the structural morphism
  \(G'\to S\); that target is used here.}%
  \endgroup

\item[\textup{(viii)}]\phantomsection\label{concl:sga3-VIB-9.2-viii}
  \emph{Let \(\mathscr P\) be a property of an \(S\)-morphism.
  Suppose that \(\mathscr P\) is local for the \textup{(fpqc)}
  topology and stable under composition. If the structural
  morphisms \(G''\to S\) and \(G'\to S\) have property
  \(\mathscr P\), then the structural morphism \(G\to S\) has
  property \(\mathscr P\) as well.}

\item[\textup{(ix)}]\phantomsection\label{concl:sga3-VIB-9.2-ix}
  \emph{If \(G\) is reduced, then \(G''\) is reduced.}

\item[\textup{(x)}]\phantomsection\label{concl:sga3-VIB-9.2-x}
  \emph{For \(G''\) to be separated over \(S\), it is necessary
  and sufficient that \(u\) \textup{(or, equivalently,
  \(\varepsilon''\))} be a closed immersion.}

\item[\textup{(xi)}]\phantomsection\label{concl:sga3-VIB-9.2-xi}
  \emph{For \(G'\) to be flat over \(S\), it is necessary and
  sufficient that \(p\) be flat \textup{(or, equivalently,
  faithfully flat)}.}

  \emph{In this case, for \(G''\) to be flat over \(S\), it is
  necessary and sufficient that \(G\) be flat over \(S\).}

\item[\textup{(xii)}]\phantomsection\label{concl:sga3-VIB-9.2-xii}
  \emph{For \(G'\) to be flat and locally of finite presentation
  over \(S\), it is necessary and sufficient that \(p:G\to G''\)
  be faithfully flat and locally of finite presentation.}

  \emph{In this case, for \(G''\) to be locally of finite
  presentation \textup{(respectively, locally of finite type, of
  finite type, smooth, \'etale, unramified, locally quasi-finite,
  quasi-finite)} over \(S\), it is enough that \(G\) have the same
  property over \(S\) \textup{(and the condition is also necessary
  in the first two cases; cf.\ \textup{(viii)})}.}

\item[\textup{(xiii)}]\phantomsection\label{concl:sga3-VIB-9.2-xiii}
  \emph{Suppose that \(G'\) is flat and of finite presentation
  over \(S\).}
  \begin{enumerate}
  \item[\textup{a)}]\phantomsection\label{concl:sga3-VIB-9.2-xiii-a}
    \emph{Then \(p\) is of finite presentation and faithfully
    flat.}
  \item[\textup{b)}]\phantomsection\label{concl:sga3-VIB-9.2-xiii-b}
    \emph{Moreover, for \(G\) to be of finite presentation over
    \(S\), it is necessary and sufficient that \(G''\) be so.}
  \end{enumerate}
\end{enumerate}

\emph{Proof.}
Recall that the equivalence relation under consideration is
universally effective (IV, \SGARefLink{sga3:iv:proposition:4:4:9}{4.4.9}). Assertions
\textup{(i)}, \textup{(iii)}, \textup{(iv)}, \textup{(v)}, and the
first assertion of \textup{(ii)} then follow, respectively, from
IV, \SGARefLink{sga3:iv:proposition:4:4:3}{4.4.3}, 5.2.2, 5.2.4, 3.4.5, and \SGARefLink{thm:sga3-VIA-3.3.2}{3.3.2}\textup{(iii)}. The second
assertion of \textup{(ii)} follows from the first, as shown by the
following cartesian diagram, since
\[
  (G\times_SG')\times_GS
\]
is isomorphic to \(G'\):
\begin{center}
  \phantomsection\label{diag:sga3-VIB-9.2-proof-cartesian}%
  {\small
  \begin{tikzcd}[row sep=large,column sep=large]
    G'
      \arrow[r,"{((\varepsilon\circ\pi'),\operatorname{id}_{G'})}"]
      \arrow[d,"\pi'"'] &
    G\times_SG'
      \arrow[r,"{\mu\circ(\operatorname{id}_G\times u)}"]
      \arrow[d,"\operatorname{pr}_1"] &
    G \arrow[d,"p"] \\
    S \arrow[r,"\varepsilon"'] &
    G \arrow[r,"p"'] &
    G''
  \end{tikzcd}}
\end{center}
Indeed, \(\varepsilon''\) is an \(S\)-section of \(G''\), hence an
immersion (EGA~I, 5.3.13); by the preceding cartesian diagram the same
is true of \(u\), which completes the proof of \textup{(ii)}.
Moreover, \textup{(vi)} is an immediate consequence of the second
cartesian diagram in \textup{(ii)}.

We prove \textup{\SGARefLink{concl:sga3-VIB-9.2-vii}{(vii)}}. By \textup{\SGARefLink{concl:sga3-VIB-9.2-i}{(i)}}, \(p\) is a covering for
the \textup{(fpqc)} topology. Hence, by \textup{(ii)}, in order to
show that \(p\) has property \(\mathscr P\), it is enough to show
that the first projection
\[
  \operatorname{pr}_1:G\times_SG'\longrightarrow G
\]
has property \(\mathscr P\). This follows from the stability of
\(\mathscr P\) under base change, since
\(\operatorname{pr}_1\) is obtained from \(\pi'\) by base change.

It is clear that \textup{(viii)} follows from \textup{(vii)}, since
\[
  \pi=\pi''\circ p,
\]
where \(\pi'':G''\to S\) denotes the structural morphism.

We prove \textup{\SGARefLink{concl:sga3-VIB-9.2-ix}{(ix)}}. By \textup{\SGARefLink{concl:sga3-VIB-9.2-i}{(i)}}, \(p\) is an epimorphism.
Since \(G\) is reduced, \(p\) factors through the immersion
\[
  G''_{\mathrm{red}}\longrightarrow G'',
\]
which is therefore also an epimorphism, and hence an isomorphism
(IV, \SGARefLink{sga3:iv:lemma:4:4:4}{4.4.4}).

We prove \textup{(x)}. If \(G''\) is separated over \(S\), then
\(\varepsilon''\) is a closed immersion by EGA~I, 5.4.6. On the
other hand, we saw in \textup{(ii)} that \(\varepsilon''\) is a
closed immersion if and only if \(u\) is. Finally, if \(u\) is a
closed immersion, then so is
\[
  \delta\circ(\operatorname{id}_G\times u):
    G\times_SG'\longrightarrow G\times_SG.
\]
It follows from
\hyperref[lem:sga3-VIB-9.2.1]{Lemma~9.2.1 below} that \(G''\) is
separated over \(S\).

Assertion \textup{(xi)} follows from \textup{(vii)} and
EGA~IV\(_2\), 2.2.13.

Assertion \textup{(xii)} follows from \textup{(vii)},
EGA~IV\(_4\), 17.7.5 and 17.7.7, and the fact that \(p\) is
universally open: for every \(s\in S\), if the underlying space of
\(G_s\) is discrete, then the same is true of the underlying space
of \(G''_s\).

Finally, assertion \textup{(xiii)} follows from \textup{(vii)},
\textup{(viii)}, and EGA~IV\(_4\), 17.7.5.
